% GENERATED FILE. Edit canonical lessons, then run npm run generate:books.
% canonical-source-hash: fnv1a64:40881d23f7379ad1
% canonical-lessons: TE-C473-dvipam, TE-C473-timingalam, TE-C473-pagadam, TE-C473-uppena, TE-C473-garjincu

\chapter{Islands and the Deep Sea}
\label{ch:te-islands-and-the-deep-sea}
\hlchaptermodality{473}

\begin{chapteropening}
\textbf{By the end of this chapter:} \emph{I can say an isle, a whale, coral, a tidal surge and to roar.}

Complete the last lesson of chapter 473: I can say an isle, a whale, coral, a tidal surge and to roar.
\end{chapteropening}

\section[\texorpdfstring{\te{ద్వీపం} (dvīpaṁ)}{ద్వీపం (dvīpaṁ)}]{\te{ద్వీపం} (dvīpaṁ) — an isle; an island}

\label{lesson:TE-C473-dvipam}



\begin{quote}
Before the new one: say the Telugu for to die out, then the Telugu for to migrate.
\end{quote}

\subsection*{You'll want to know: \te{ద్వీపం}}
\textbf{\te{ద్వీపం}} — \emph{dvīpaṁ} — \textquotedblleft{}an isle; an island\textquotedblright{}.

\begin{quote}
\te{ఆ} \te{ద్వీపానికి} \te{పడవలో} \te{వెళ్ళాలి}.
\end{quote}

\begin{quote}
\emph{ā dvīpāniki paḍavalō veḷḷāli.}
\end{quote}

\begin{quote}
You have to go to that island by boat.
\end{quote}

A formal word; \te{దీవి} is the everyday one.

\begin{cousinweb}
\textbf{\te{ద్వీపం}} (\emph{dvīpaṁ}) is borrowed from Sanskrit \emph{dvīpa}, \textquotedblleft{}land, an island\textquotedblright{} (CDIAL 6691), which CDIAL takes apart as \textquotedblleft{}with water on two sides\textquotedblright{}.
\end{cousinweb}

\begin{grammarlens}[title={What you've built}]
One more word for getting around.
\end{grammarlens}

\subsection*{Your turn}
\begin{itemize}
\raggedright
  \item \emph{Say it:} \emph{dvīpaṁ}
  \item \emph{Say it:} \emph{dvīpaṁ}, once more
  \item \emph{Say it:} \emph{ā dvīpāniki paḍavalō veḷḷāli.} — You have to go to that island by boat.
  \item \emph{Recall:} say the Telugu for to die out, then the Telugu for to migrate, then say \emph{dvīpaṁ} again
\end{itemize}

\begin{tcolorbox}[breakable,colback=teal!4,colframe=teal!35!black,title={Before you move on}]
What does \te{ద్వీపం} mean? (\textquotedblleft{}An isle; an island\textquotedblright{}.) Say it once more.
\end{tcolorbox}
\section[\texorpdfstring{\te{తిమింగలం} (timiṅgalaṁ)}{తిమింగలం (timiṅgalaṁ)}]{\te{తిమింగలం} (timiṅgalaṁ) — a whale}

\label{lesson:TE-C473-timingalam}



\begin{quote}
Before the new one: say the Telugu for to migrate, then the Telugu for an isle.
\end{quote}

\subsection*{You'll want to know: \te{తిమింగలం}}
\textbf{\te{తిమింగలం}} — \emph{timiṅgalaṁ} — \textquotedblleft{}a whale\textquotedblright{}.

\begin{quote}
\te{సముద్రంలో} \te{పెద్ద} \te{తిమింగలం} \te{కనిపించింది}.
\end{quote}

\begin{quote}
\emph{samudraṁlō pedda timiṅgalaṁ kanipiñcindi.}
\end{quote}

\begin{quote}
A huge whale appeared in the sea.
\end{quote}

The biggest animal in the sea.

\begin{cousinweb}
\textbf{\te{తిమింగలం}} (\emph{timiṅgalaṁ}) is borrowed from Sanskrit \emph{timiṅgila}, the name of a huge sea creature, literally \textquotedblleft{}swallower of the \emph{timi} fish\textquotedblright{} (Monier-Williams).
\end{cousinweb}

\begin{grammarlens}[title={What you've built}]
One more word for getting around.
\end{grammarlens}

\subsection*{Your turn}
\begin{itemize}
\raggedright
  \item \emph{Say it:} \emph{timiṅgalaṁ}
  \item \emph{Say it:} \emph{timiṅgalaṁ}, once more
  \item \emph{Say it:} \emph{samudraṁlō pedda timiṅgalaṁ kanipiñcindi.} — A huge whale appeared in the sea.
  \item \emph{Recall:} say the Telugu for to migrate, then the Telugu for an isle, then say \emph{timiṅgalaṁ} again
\end{itemize}

\begin{tcolorbox}[breakable,colback=teal!4,colframe=teal!35!black,title={Before you move on}]
What does \te{తిమింగలం} mean? (\textquotedblleft{}A whale\textquotedblright{}.) Say it once more.
\end{tcolorbox}
\section[\texorpdfstring{\te{పగడం} (pagaḍaṁ)}{పగడం (pagaḍaṁ)}]{\te{పగడం} (pagaḍaṁ) — coral}

\label{lesson:TE-C473-pagadam}



\begin{quote}
Before the new one: say the Telugu for an isle, then the Telugu for a whale.
\end{quote}

\subsection*{You'll want to know: \te{పగడం}}
\textbf{\te{పగడం}} — \emph{pagaḍaṁ} — \textquotedblleft{}coral\textquotedblright{}.

\begin{quote}
\te{ఆ} \te{ద్వీపం} \te{చుట్టూ} \te{పగడాల} \te{దిబ్బలు} \te{ఉన్నాయి}.
\end{quote}

\begin{quote}
\emph{ā dvīpaṁ cuṭṭū pagaḍāla dibbalu unnāyi.}
\end{quote}

\begin{quote}
There are coral reefs round that island.
\end{quote}

\te{పగడాల} \te{దిబ్బ}, a coral reef. Red coral is set in jewellery too.

\begin{cousinweb}
\textbf{\te{పగడం}} (\emph{pagaḍaṁ}) is listed in DEDR as Telugu \emph{pagaḍamu}, \textquotedblleft{}red coral\textquotedblright{} (DEDR 3998), beside Tamil \emph{pavaḷam}; DEDR compares Sanskrit \emph{pravāla}, and which side borrowed is not settled.
\end{cousinweb}

\begin{grammarlens}[title={What you've built}]
One more word for getting around.
\end{grammarlens}

\subsection*{Your turn}
\begin{itemize}
\raggedright
  \item \emph{Say it:} \emph{pagaḍaṁ}
  \item \emph{Say it:} \emph{pagaḍaṁ}, once more
  \item \emph{Say it:} \emph{ā dvīpaṁ cuṭṭū pagaḍāla dibbalu unnāyi.} — There are coral reefs round that island.
  \item \emph{Recall:} say the Telugu for an isle, then the Telugu for a whale, then say \emph{pagaḍaṁ} again
\end{itemize}

\begin{tcolorbox}[breakable,colback=teal!4,colframe=teal!35!black,title={Before you move on}]
What does \te{పగడం} mean? (\textquotedblleft{}Coral\textquotedblright{}.) Say it once more.
\end{tcolorbox}
\section[\texorpdfstring{\te{ఉప్పెన} (uppena)}{ఉప్పెన (uppena)}]{\te{ఉప్పెన} (uppena) — a tidal surge}

\label{lesson:TE-C473-uppena}



\begin{quote}
Before the new one: say the Telugu for a whale, then the Telugu for coral.
\end{quote}

\subsection*{You'll want to know: \te{ఉప్పెన}}
\textbf{\te{ఉప్పెన}} — \emph{uppena} — \textquotedblleft{}a tidal surge\textquotedblright{}.

\begin{quote}
\te{తుపాను} \te{వల్ల} \te{తీరంలో} \te{ఉప్పెన} \te{వచ్చింది}.
\end{quote}

\begin{quote}
\emph{tupānu valla tīraṁlō uppena vaccindi.}
\end{quote}

\begin{quote}
Because of the storm a tidal surge hit the coast.
\end{quote}

The sea rising over the land in a storm.

\begin{cousinweb}
\textbf{\te{ఉప్పెన}} (\emph{uppena}) is native Dravidian: DEDR 666 lists Telugu \emph{uppena}, \textquotedblleft{}a swell of the sea, an inundation\textquotedblright{}, beside \emph{ubbu}, \textquotedblleft{}to swell, to overflow\textquotedblright{}. The sea swells over the land.
\end{cousinweb}

\begin{grammarlens}[title={What you've built}]
One more word for getting around.
\end{grammarlens}

\subsection*{Your turn}
\begin{itemize}
\raggedright
  \item \emph{Say it:} \emph{uppena}
  \item \emph{Say it:} \emph{uppena}, once more
  \item \emph{Say it:} \emph{tupānu valla tīraṁlō uppena vaccindi.} — Because of the storm a tidal surge hit the coast.
  \item \emph{Recall:} say the Telugu for a whale, then the Telugu for coral, then say \emph{uppena} again
\end{itemize}

\begin{tcolorbox}[breakable,colback=teal!4,colframe=teal!35!black,title={Before you move on}]
What does \te{ఉప్పెన} mean? (\textquotedblleft{}A tidal surge\textquotedblright{}.) Say it once more.
\end{tcolorbox}
\section[\texorpdfstring{\te{గర్జించు} (garjiñcu)}{గర్జించు (garjiñcu)}]{\te{గర్జించు} (garjiñcu) — to roar}

\label{lesson:TE-C473-garjincu}



\begin{quote}
Before the new one: say the Telugu for coral, then the Telugu for a tidal surge.
\end{quote}

\subsection*{You'll want to know: \te{గర్జించు}}
\textbf{\te{గర్జించు}} — \emph{garjiñcu} — \textquotedblleft{}to roar\textquotedblright{}.

\begin{quote}
\te{రాత్రంతా} \te{సముద్రం} \te{గర్జించింది}.
\end{quote}

\begin{quote}
\emph{rātrantā samudraṁ garjiñcindi.}
\end{quote}

\begin{quote}
The sea roared all night.
\end{quote}

Past (I …): \te{గర్జించాను}, \emph{garjiñcānu}.

Said of a lion, of thunder and of an angry sea.

\begin{cousinweb}
\textbf{\te{గర్జించు}} (\emph{garjiñcu}) is built from Sanskrit \emph{garj}, \textquotedblleft{}to thunder, to roar\textquotedblright{} (CDIAL 4046, \emph{garjati}), with the Telugu verb ending -\te{ఇంచు}.
\end{cousinweb}

\begin{grammarlens}[title={What you've built}]
One more word for getting around.
\end{grammarlens}

\subsection*{Your turn}
\begin{itemize}
\raggedright
  \item \emph{Say it:} \emph{garjiñcu}
  \item \emph{Say it:} \emph{garjiñcu}, once more
  \item \emph{Say it:} \emph{rātrantā samudraṁ garjiñcindi.} — The sea roared all night.
  \item \emph{Recall:} say the Telugu for coral, then the Telugu for a tidal surge, then say \emph{garjiñcu} again
\end{itemize}

\begin{tcolorbox}[breakable,colback=teal!4,colframe=teal!35!black,title={Before you move on}]
What does \te{గర్జించు} mean? (\textquotedblleft{}To roar\textquotedblright{}.) Say it once more.
\end{tcolorbox}
